% ============================================
% Intelligent Document Categorization
% NLP and Text Analysis - Project Report
% ============================================
\documentclass[12pt, a4paper]{article}

% --- Packages ---
\usepackage[utf8]{inputenc}
\usepackage[T1]{fontenc}
\usepackage{lmodern}
\usepackage[margin=1in]{geometry}
\usepackage{graphicx}
\usepackage{booktabs}
\usepackage{tabularx}
\usepackage{longtable}
\usepackage{amsmath}
\usepackage{amssymb}
\usepackage{hyperref}
\usepackage{xcolor}
\usepackage{listings}
\usepackage{caption}
\usepackage{subcaption}
\usepackage{float}
\usepackage{enumitem}
\usepackage{titlesec}
\usepackage{fancyhdr}
\usepackage{setspace}
\usepackage{parskip}
\usepackage{cite}

% --- Hyperlink Setup ---
\hypersetup{
    colorlinks=true,
    linkcolor=blue!70!black,
    citecolor=green!50!black,
    urlcolor=blue!70!black,
}

% --- Code Listing Style ---
\lstdefinestyle{python}{
    language=Python,
    basicstyle=\small\ttfamily,
    keywordstyle=\color{blue!70!black}\bfseries,
    stringstyle=\color{red!60!black},
    commentstyle=\color{green!50!black}\itshape,
    numberstyle=\tiny\color{gray},
    numbers=left,
    numbersep=5pt,
    frame=single,
    breaklines=true,
    breakatwhitespace=true,
    tabsize=4,
    showstringspaces=false,
    backgroundcolor=\color{gray!5},
    rulecolor=\color{gray!30},
}
\lstset{style=python}

% --- Header/Footer ---
\pagestyle{fancy}
\fancyhf{}
\fancyhead[L]{\small NLP and Text Analysis}
\fancyhead[R]{\small Document Categorization}
\fancyfoot[C]{\thepage}
\renewcommand{\headrulewidth}{0.4pt}
\renewcommand{\footrulewidth}{0.4pt}

% --- Line Spacing ---
\onehalfspacing

% ============================================
% TITLE PAGE
% ============================================
\begin{document}

\begin{titlepage}
    \centering
    \vspace*{1cm}

    {\LARGE\bfseries Intelligent Document Categorization\\using Hugging Face Transformers\\with Performance Benchmarking\par}

    \vspace{1.5cm}

    {\large\textit{A Project Report}\par}
    {\large\textit{for the course}\par}
    {\large\textbf{Natural Language Processing and Text Analysis}\par}

    \vspace{2cm}

    {\large\textbf{Submitted by:}\par}
    \vspace{0.5cm}
    \begin{tabular}{ll}
        \toprule
        \textbf{Name} & \textbf{Roll Number} \\
        \midrule
        Member 1 Name & Roll No. 1 \\
        Member 2 Name & Roll No. 2 \\
        \bottomrule
    \end{tabular}

    \vspace{2cm}

    {\large\textbf{Under the Guidance of:}\par}
    {\large Professor Name\par}
    {\large Department Name\par}

    \vfill

    {\large Institution/University Name\par}
    \vspace{0.5cm}
    {\large Academic Year 2025--2026\par}
\end{titlepage}

% ============================================
% TABLE OF CONTENTS
% ============================================
\newpage
\tableofcontents
\newpage
\listoffigures
\listoftables
\newpage

% ============================================
% 1. INTRODUCTION
% ============================================
\section{Introduction}

\subsection{Background}
Document categorization, also known as text classification, is one of the fundamental tasks in Natural Language Processing (NLP). It involves automatically assigning predefined category labels to text documents based on their content. This task has numerous real-world applications including email filtering, news article organization, sentiment analysis, spam detection, and content recommendation systems.

With the advent of deep learning and particularly Transformer-based architectures, the landscape of text classification has undergone a significant transformation. Pre-trained language models such as BERT \cite{devlin2019bert}, DistilBERT \cite{sanh2019distilbert}, and RoBERTa \cite{liu2019roberta} have achieved state-of-the-art performance across a wide range of NLP benchmarks by leveraging transfer learning from large-scale unlabeled corpora.

\subsection{Problem Statement}
This project addresses the task of multi-class document categorization by implementing and comparing multiple approaches:
\begin{enumerate}
    \item \textbf{Traditional ML baselines} using TF-IDF feature extraction with Logistic Regression, Naive Bayes, and Support Vector Machines.
    \item \textbf{Fine-tuned Transformer models} including BERT, DistilBERT, and RoBERTa using the Hugging Face \texttt{transformers} library.
    \item \textbf{Zero-shot classification} using BART-large-MNLI to evaluate performance without any task-specific training.
\end{enumerate}

A key contribution of this work is the comprehensive \textbf{performance benchmarking} that goes beyond accuracy to include macro/weighted F1 scores, inference latency, model size, and computational requirements.

\subsection{Objectives}
\begin{itemize}
    \item Implement an end-to-end document categorization pipeline from data preparation to deployment.
    \item Compare traditional ML approaches with modern Transformer-based methods.
    \item Evaluate the effectiveness of zero-shot classification as a training-free alternative.
    \item Conduct comprehensive benchmarking across multiple performance dimensions.
    \item Develop a live demo application for interactive document categorization.
\end{itemize}

\subsection{Scope and Limitations}
The project focuses on English-language news article classification using two standard benchmark datasets. Experiments are conducted on an NVIDIA RTX 2050 GPU with 4GB VRAM, which constrains model sizes and batch configurations. All Transformer models use the \texttt{base} variant (approximately 110M parameters) due to hardware limitations.

% ============================================
% 2. LITERATURE REVIEW
% ============================================
\section{Literature Review}

\subsection{Traditional Text Classification}
Traditional text classification methods rely on handcrafted feature extraction followed by machine learning classifiers. The \textbf{Bag-of-Words (BoW)} and \textbf{TF-IDF (Term Frequency-Inverse Document Frequency)} representations have been the dominant feature extraction methods for decades \cite{manning2008introduction}.

\textbf{TF-IDF} weighs terms by their frequency within a document (TF) relative to their frequency across the corpus (IDF):
\begin{equation}
    \text{TF-IDF}(t, d) = \text{TF}(t, d) \times \log\left(\frac{N}{\text{DF}(t)}\right)
\end{equation}
where $t$ is a term, $d$ is a document, $N$ is the total number of documents, and $\text{DF}(t)$ is the number of documents containing term $t$.

Common classifiers used with TF-IDF include:
\begin{itemize}
    \item \textbf{Logistic Regression:} A linear model that uses the softmax function for multi-class classification.
    \item \textbf{Naive Bayes:} A probabilistic classifier based on Bayes' theorem with the ``naive'' assumption of feature independence.
    \item \textbf{Support Vector Machines:} Maximum-margin classifiers that find optimal hyperplanes separating classes.
\end{itemize}

\subsection{Transformer Architecture}
The Transformer architecture, introduced by Vaswani et al. \cite{vaswani2017attention}, revolutionized NLP through its self-attention mechanism. The core \textbf{scaled dot-product attention} is computed as:
\begin{equation}
    \text{Attention}(Q, K, V) = \text{softmax}\left(\frac{QK^T}{\sqrt{d_k}}\right)V
\end{equation}
where $Q$, $K$, $V$ are the query, key, and value matrices, and $d_k$ is the dimension of the key vectors.

\subsection{BERT and Variants}
\textbf{BERT} (Bidirectional Encoder Representations from Transformers) \cite{devlin2019bert} introduced bidirectional pre-training using Masked Language Modeling (MLM) and Next Sentence Prediction (NSP). It processes text in both directions simultaneously, capturing richer contextual representations.

\textbf{DistilBERT} \cite{sanh2019distilbert} is a distilled version of BERT that retains 97\% of BERT's performance while being 60\% faster and 40\% smaller, achieved through knowledge distillation during pre-training.

\textbf{RoBERTa} \cite{liu2019roberta} optimizes BERT's pre-training by removing NSP, using dynamic masking, larger batch sizes, and more training data.

\subsection{Zero-Shot Classification}
Zero-shot classification enables categorizing text into arbitrary labels without task-specific training. BART-large-MNLI \cite{lewis2020bart} reformulates classification as a natural language inference (NLI) task, checking if the document ``entails'' each candidate label.

% ============================================
% 3. METHODOLOGY
% ============================================
\section{Methodology}

\subsection{System Architecture}
Figure~\ref{fig:architecture} illustrates the overall system architecture of the document categorization pipeline.

\begin{figure}[H]
    \centering
    \fbox{\parbox{0.9\textwidth}{\centering
        \textit{[Insert system architecture diagram from results/figures/ here]}\\
        \vspace{0.5cm}
        Raw Datasets $\rightarrow$ Data Preparation $\rightarrow$ EDA $\rightarrow$ Model Training $\rightarrow$ Benchmarking $\rightarrow$ Demo App
    }}
    \caption{System architecture of the document categorization pipeline.}
    \label{fig:architecture}
\end{figure}

\subsection{Datasets}
Two standard benchmark datasets are used for evaluation:

\begin{table}[H]
    \centering
    \caption{Dataset summary.}
    \label{tab:datasets}
    \begin{tabular}{lccl}
        \toprule
        \textbf{Dataset} & \textbf{Samples} & \textbf{Classes} & \textbf{Categories} \\
        \midrule
        AG News & 127,600 & 4 & World, Sports, Business, Sci/Tech \\
        \bottomrule
    \end{tabular}
\end{table}

\textbf{AG News} (accessible via Hugging Face at \url{https://huggingface.co/datasets/fancyzhx/ag_news}) provides a large-scale benchmark with 120K training samples across four balanced news categories. A subset of 20,000 training samples is used for computational feasibility on the available hardware.

\subsection{Data Preprocessing}
Text preprocessing includes the following steps:
\begin{enumerate}
    \item Conversion to lowercase.
    \item Removal of special characters (retaining only alphanumeric characters and spaces).
    \item Stripping of extra whitespace.
    \item Stratified train/validation/test splitting (80/10/10 for BBC News; built-in splits for AG News with a 20K training subset for computational feasibility).
\end{enumerate}

\subsection{Feature Extraction: TF-IDF}
For baseline models, TF-IDF vectorization is applied with the following configuration:
\begin{itemize}
    \item Maximum features: 50,000
    \item N-gram range: (1, 2) — unigrams and bigrams
    \item Sublinear TF scaling
\end{itemize}

\subsection{Baseline Models}
Three traditional ML classifiers are trained on TF-IDF features:
\begin{enumerate}
    \item \textbf{Logistic Regression:} Multinomial classification with L2 regularization ($C=1.0$).
    \item \textbf{Multinomial Naive Bayes:} Laplace smoothing ($\alpha=1.0$).
    \item \textbf{Linear SVM (LinearSVC):} Maximum margin classifier ($C=1.0$).
\end{enumerate}

\subsection{Transformer Fine-Tuning}
Pre-trained Transformer models are fine-tuned for sequence classification using the Hugging Face \texttt{Trainer} API. Table~\ref{tab:hyperparams} lists the training hyperparameters, optimized for a 4GB VRAM constraint.

\begin{table}[H]
    \centering
    \caption{Transformer fine-tuning hyperparameters.}
    \label{tab:hyperparams}
    \begin{tabular}{ll}
        \toprule
        \textbf{Parameter} & \textbf{Value} \\
        \midrule
        Max Sequence Length & 128 tokens \\
        Training Batch Size & 8 \\
        Gradient Accumulation Steps & 4 (effective batch = 32) \\
        Learning Rate & $2 \times 10^{-5}$ \\
        Weight Decay & 0.01 \\
        Epochs & 3 \\
        Warmup Ratio & 0.1 \\
        Mixed Precision (FP16) & Enabled \\
        Optimizer & AdamW \\
        Early Stopping Patience & 2 epochs \\
        \bottomrule
    \end{tabular}
\end{table}

\subsection{Zero-Shot Classification}
The \texttt{facebook/bart-large-mnli} model is used for zero-shot classification. Given a document and a set of candidate labels, the model predicts the most likely label without any task-specific fine-tuning. Evaluation is performed on a subset of 500 test samples due to the higher computational cost of this approach.

\subsection{Evaluation Metrics}
The following metrics are computed for comprehensive evaluation:
\begin{itemize}
    \item \textbf{Accuracy:} $\frac{\text{Correct Predictions}}{\text{Total Predictions}}$
    \item \textbf{Macro F1:} Unweighted mean of per-class F1 scores:
    \begin{equation}
        F1_{\text{macro}} = \frac{1}{C}\sum_{i=1}^{C} \frac{2 \cdot P_i \cdot R_i}{P_i + R_i}
    \end{equation}
    \item \textbf{Weighted F1:} F1 weighted by class support.
    \item \textbf{Inference Latency:} Average time per prediction (milliseconds).
    \item \textbf{Model Size:} Total parameters and disk footprint.
\end{itemize}

% ============================================
% 4. EXPERIMENTAL RESULTS
% ============================================
\section{Experimental Results}

\subsection{Exploratory Data Analysis}
% Include EDA figures here
\textit{[Insert class distribution, text length histograms, and word clouds from results/figures/ after running eda.py]}

\subsection{Baseline Model Results}
\begin{table}[H]
    \centering
    \caption{Baseline model performance on AG News dataset.}
    \label{tab:baseline_ag}
    \begin{tabular}{lccc}
        \toprule
        \textbf{Model} & \textbf{Accuracy} & \textbf{Macro F1} & \textbf{Latency (ms)} \\
        \midrule
        Logistic Regression & 0.8932 & 0.8928 & 0.40 \\
        Naive Bayes & 0.8909 & 0.8905 & 1.13 \\
        SVM (LinearSVC) & \textbf{0.9011} & \textbf{0.9008} & 0.42 \\
        \bottomrule
    \end{tabular}
\end{table}

Among the traditional baselines, SVM (LinearSVC) achieved the highest accuracy of 90.11\% with a Macro F1 of 0.9008. All three baselines demonstrated sub-millisecond to low-millisecond inference latency, making them highly suitable for real-time applications. Logistic Regression was the fastest at 0.40 ms per sample.


\subsection{Transformer Model Results}
\begin{table}[H]
    \centering
    \caption{Transformer model performance on AG News dataset.}
    \label{tab:transformer_ag}
    \begin{tabular}{lcccc}
        \toprule
        \textbf{Model} & \textbf{Accuracy} & \textbf{Macro F1} & \textbf{Params (M)} & \textbf{Latency (ms)} \\
        \midrule
        BERT-base & \textbf{0.9229} & \textbf{0.9227} & 110 & 34.96 \\
        DistilBERT & 0.9228 & 0.9227 & 66 & \textbf{12.55} \\
        RoBERTa-base & 0.9205 & 0.9204 & 125 & 19.46 \\
        BART Zero-Shot & 0.7140 & 0.6961 & 407 & 186.72 \\
        \bottomrule
    \end{tabular}
\end{table}

BERT and DistilBERT achieved nearly identical accuracy (92.29\% vs 92.28\%), while DistilBERT was 2.8$\times$ faster in inference (12.55 ms vs 34.96 ms) with only 60\% of the parameters. RoBERTa performed slightly lower at 92.05\%. The zero-shot BART model, requiring no task-specific training, achieved 71.40\% accuracy --- a 21\% gap compared to fine-tuned models.

\subsection{Comparative Analysis}

Table~\ref{tab:comparison} presents the unified comparison across all seven models.

\begin{table}[H]
    \centering
    \caption{Unified performance comparison across all models on AG News.}
    \label{tab:comparison}
    \begin{tabular}{lcccc}
        \toprule
        \textbf{Model} & \textbf{Accuracy} & \textbf{Macro F1} & \textbf{Latency (ms)} & \textbf{Type} \\
        \midrule
        BERT-base & \textbf{0.9229} & \textbf{0.9227} & 34.96 & Transformer \\
        DistilBERT & 0.9228 & 0.9227 & 12.55 & Transformer \\
        RoBERTa-base & 0.9205 & 0.9204 & 19.46 & Transformer \\
        SVM (LinearSVC) & 0.9011 & 0.9008 & \textbf{0.42} & Baseline \\
        Logistic Regression & 0.8932 & 0.8928 & 0.40 & Baseline \\
        Naive Bayes & 0.8909 & 0.8905 & 1.13 & Baseline \\
        BART Zero-Shot & 0.7140 & 0.6961 & 186.72 & Zero-Shot \\
        \bottomrule
    \end{tabular}
\end{table}

\subsection{Receiver Operating Characteristic (ROC) and Multi-Class AUC Analysis}
To evaluate classification discriminability independent of arbitrary decision thresholds, One-vs-Rest (OvR) Receiver Operating Characteristic (ROC) curves and Area Under the Curve (AUC) metrics were computed across all probability-producing models on the 7,600-sample AG News test set.

\begin{table}[H]
    \centering
    \caption{Multi-class ROC-AUC scores (One-vs-Rest) across benchmarked models on AG News.}
    \label{tab:roc_auc}
    \begin{tabular}{lcccccc}
        \toprule
        \textbf{Model} & \textbf{Macro AUC} & \textbf{Micro AUC} & \textbf{World} & \textbf{Sports} & \textbf{Business} & \textbf{Sci/Tech} \\
        \midrule
        BERT-base & \textbf{0.9868} & \textbf{0.9889} & \textbf{0.9839} & \textbf{0.9974} & \textbf{0.9812} & \textbf{0.9843} \\
        DistilBERT & 0.9861 & 0.9883 & 0.9826 & 0.9973 & 0.9801 & 0.9839 \\
        RoBERTa-base & 0.9852 & 0.9869 & 0.9828 & 0.9964 & 0.9789 & 0.9821 \\
        SVM (LinearSVC) & 0.9774 & 0.9786 & 0.9745 & 0.9944 & 0.9662 & 0.9739 \\
        Logistic Regression & 0.9762 & 0.9778 & 0.9715 & 0.9935 & 0.9661 & 0.9733 \\
        Naive Bayes & 0.9738 & 0.9756 & 0.9654 & 0.9939 & 0.9634 & 0.9719 \\
        \bottomrule
    \end{tabular}
\end{table}

Across all models, the \textit{Sports} class achieved the highest AUC ($>0.993$), reflecting a distinct domain-specific vocabulary. Conversely, \textit{Business} and \textit{Sci/Tech} exhibited the lowest AUCs ($0.963$--$0.981$) due to shared technical and commercial terms in technology corporate news. DistilBERT closely matches BERT with a Macro AUC of 0.9861, demonstrating excellent probability calibration.

\subsection{Key Findings}
\begin{enumerate}
    \item \textbf{Transformer models outperform traditional baselines} by approximately 2 percentage points (92.3\% vs 90.1\%), demonstrating the value of pre-trained contextual representations over TF-IDF features.
    \item \textbf{DistilBERT offers the best accuracy-efficiency trade-off}, matching BERT's accuracy (92.28\% vs 92.29\%) while being 2.8$\times$ faster and using only 60\% of the parameters (66M vs 110M).
    \item \textbf{Zero-shot classification lags significantly} at 71.4\% accuracy, a 21\% gap compared to fine-tuned models, confirming that task-specific fine-tuning remains essential for high-accuracy applications.
    \item \textbf{Inference latency varies by two orders of magnitude}: TF-IDF baselines operate at 0.4 ms per sample while Transformer models require 12--35 ms and zero-shot requires 187 ms, a critical consideration for deployment.
    \item \textbf{SVM is the strongest traditional baseline} at 90.1\%, outperforming Logistic Regression and Naive Bayes while maintaining sub-millisecond latency.
\end{enumerate}

% ============================================
% 5. DISCUSSION
% ============================================
\section{Discussion}

\subsection{Impact of Training Data Size}
The AG News training subset was limited to 20,000 samples (from the full 120,000) due to GPU memory constraints. Despite using only 16.7\% of the available training data, Transformer models achieved over 92\% accuracy, demonstrating the strong transfer learning capabilities of pre-trained language models. Training on the full dataset would likely yield further improvements.

\subsection{Computational Considerations}
Training was conducted on an NVIDIA RTX 2050 GPU with 4GB VRAM. To accommodate the memory constraints:
\begin{itemize}
    \item Batch size was reduced to 8 with gradient accumulation over 4 steps.
    \item Mixed precision training (FP16) was enabled.
    \item Maximum sequence length was limited to 128 tokens.
    \item AG News training set was subsampled to 20,000 examples.
\end{itemize}

These constraints may slightly underestimate the full potential of Transformer models compared to training with larger batch sizes and full datasets.

\subsection{Zero-Shot vs. Fine-Tuned Models}
The zero-shot approach using BART-large-MNLI offers significant practical advantages --- no training data, no fine-tuning, and immediate deployment capability. However, the 21 percentage point accuracy gap (71.4\% vs 92.3\%) demonstrates that task-specific fine-tuning remains essential when labeled data is available and high accuracy is required.

\subsection{Practical Recommendations}
Based on the experimental results:
\begin{itemize}
    \item For \textbf{high-accuracy applications}: Use fine-tuned BERT or DistilBERT (92.3\% accuracy).
    \item For \textbf{resource-constrained deployment}: Use DistilBERT (66M params, 12.5 ms latency) or TF-IDF + SVM (0.42 ms latency).
    \item For \textbf{rapid prototyping without training data}: Use zero-shot classification (71.4\% accuracy, no training required).
    \item For \textbf{real-time, high-throughput applications}: Use TF-IDF baselines (sub-millisecond latency, 90\% accuracy).
\end{itemize}

% ============================================
% 6. CONCLUSION
% ============================================
\section{Conclusion}

This project successfully designed, implemented, and benchmarked an end-to-end intelligent document categorization framework on the AG News corpus. By systematically comparing traditional TF-IDF classifiers against state-of-the-art Hugging Face Transformers and zero-shot NLI models across accuracy, macro/micro F1, multi-class ROC-AUC, and inference latency, this work provides clear empirical guidelines for deploying NLP systems under diverse computational constraints.

\subsection{Key Findings}
\begin{itemize}
    \item \textbf{DistilBERT Delivers the Optimal Trade-Off:} DistilBERT achieved 92.28\% accuracy (essentially matching BERT's 92.29\%) and a 0.9861 Macro ROC-AUC while operating 2.8$\times$ faster with 40\% fewer parameters.
    \item \textbf{Linear SVM is Ideal for Ultra-Low Latency:} TF-IDF + LinearSVC provided 90.11\% accuracy at 0.42 ms latency, delivering 97.6\% of BERT's accuracy at less than 1/80th the computational cost.
    \item \textbf{Zero-Shot Viability in Cold-Start Scenarios:} BART-large-MNLI attained 71.40\% accuracy without requiring any training examples or parameter updates.
    \item \textbf{Class Separability Characteristics:} The \textit{Sports} category exhibited the cleanest separation ($>0.99$ AUC), whereas \textit{Business} and \textit{Sci/Tech} exhibited the highest semantic cross-contamination due to corporate tech reporting.
\end{itemize}

\subsection{Challenges Faced}
\begin{itemize}
    \item \textbf{Hardware and Memory Constraints (4GB VRAM):} Training 110M+ parameter Transformers on a mobile workstation GPU (NVIDIA RTX 2050 with 4GB VRAM) led to CUDA Out-Of-Memory exceptions under standard batch sizes. This was successfully resolved by introducing mixed precision (FP16), small per-device batch sizes (8), gradient accumulation over 4 steps, and sequence truncation to 128 tokens.
    \item \textbf{Multi-Class Probability Calibration:} Margin-based linear classifiers (LinearSVC) do not natively produce calibrated posterior probabilities required for ROC analysis. A softmax-calibrated decision function was formulated to enable unbiased One-vs-Rest ROC curve generation.
    \item \textbf{Latency-Accuracy Trade-off:} Balancing deep contextual representations against real-time throughput constraints required systematic profiling of CPU and GPU inference latencies.
\end{itemize}

\subsection{Future Scope}
\begin{itemize}
    \item \textbf{Parameter-Efficient Fine-Tuning (PEFT):} Integrating LoRA and QLoRA to enable fine-tuning of modern open-weights instruction models (such as Llama 3 and Mistral 7B) within consumer hardware limits.
    \item \textbf{Multilingual Document Categorization:} Expanding the pipeline using Multilingual BERT (mBERT) and XLM-RoBERTa for cross-lingual enterprise document indexing.
    \item \textbf{Hierarchical and Long-Document Attention:} Incorporating chunking and hierarchical pooling mechanisms or Longformer architectures to categorize multi-page contracts and legal briefs without truncation.
    \item \textbf{Tiered Inference Ensembles:} Deploying a dual-stage pipeline where high-confidence inputs are resolved by Linear SVM (0.4 ms) and ambiguous cases are escalated to DistilBERT.
\end{itemize}

\subsection{Applications of the Project}
\begin{itemize}
    \item \textbf{Automated Newsroom and Feed Aggregation:} Instant routing and tagging of thousands of syndicated news stories and RSS feeds.
    \item \textbf{Enterprise Customer Support Triage:} Automatically directing incoming customer inquiries, bug reports, and billing tickets to the appropriate specialized teams.
    \item \textbf{Legal and Regulatory Compliance Filing:} Rapid indexing, discovery, and archiving of regulatory disclosures and contracts.
    \item \textbf{Brand Safety and Forum Moderation:} Categorizing user-generated publications and posts to prevent inappropriate placement and policy infractions.
\end{itemize}

% ============================================
% REFERENCES
% ============================================
\newpage
\bibliographystyle{plain}

\begin{thebibliography}{10}

\bibitem{devlin2019bert}
J.~Devlin, M.-W. Chang, K.~Lee, and K.~Toutanova.
\newblock BERT: Pre-training of Deep Bidirectional Transformers for Language Understanding.
\newblock In \textit{Proceedings of NAACL-HLT}, pages 4171--4186, 2019.

\bibitem{sanh2019distilbert}
V.~Sanh, L.~Debut, J.~Chaumond, and T.~Wolf.
\newblock DistilBERT, a distilled version of BERT: smaller, faster, cheaper and lighter.
\newblock In \textit{NeurIPS Workshop on Energy Efficient Machine Learning and Cognitive Computing}, 2019.

\bibitem{liu2019roberta}
Y.~Liu, M.~Ott, N.~Goyal, J.~Du, M.~Joshi, D.~Chen, O.~Levy, M.~Lewis, L.~Zettlemoyer, and V.~Stoyanov.
\newblock RoBERTa: A Robustly Optimized BERT Pretraining Approach.
\newblock \textit{arXiv preprint arXiv:1907.11692}, 2019.

\bibitem{lewis2020bart}
M.~Lewis, Y.~Liu, N.~Goyal, M.~Ghazvininejad, A.~Mohamed, O.~Levy, V.~Stoyanov, and L.~Zettlemoyer.
\newblock BART: Denoising Sequence-to-Sequence Pre-training for Natural Language Generation, Translation, and Comprehension.
\newblock In \textit{Proceedings of ACL}, pages 7871--7880, 2020.

\bibitem{vaswani2017attention}
A.~Vaswani, N.~Shazeer, N.~Parmar, J.~Uszkoreit, L.~Jones, A.~N. Gomez, \L.~Kaiser, and I.~Polosukhin.
\newblock Attention Is All You Need.
\newblock In \textit{Advances in Neural Information Processing Systems (NeurIPS)}, pages 5998--6008, 2017.

\bibitem{manning2008introduction}
C.~D. Manning, P.~Raghavan, and H.~Sch\"utze.
\newblock \textit{Introduction to Information Retrieval}.
\newblock Cambridge University Press, 2008.

\bibitem{wolf2020transformers}
T.~Wolf, L.~Debut, V.~Sanh, J.~Chaumond, C.~Delangue, A.~Moi, P.~Cistac, T.~Rault, R.~Louf, M.~Funtowicz, J.~Davison, S.~Shleifer, P.~von Platen, C.~Ma, Y.~Jernite, J.~Plu, C.~Xu, T.~Le Scao, S.~Gugger, M.~Drame, Q.~Lhoest, and A.~Rush.
\newblock Transformers: State-of-the-Art Natural Language Processing.
\newblock In \textit{Proceedings of EMNLP: System Demonstrations}, pages 38--45, 2020.

\bibitem{yin2019benchmarking}
W.~Yin, J.~Hay, and D.~Roth.
\newblock Benchmarking Zero-shot Text Classification: Datasets, Evaluation and Entailment Approach.
\newblock In \textit{Proceedings of EMNLP-IJCNLP}, pages 3914--3923, 2019.

\bibitem{zhang2015character}
X.~Zhang, J.~Zhao, and Y.~LeCun.
\newblock Character-level Convolutional Networks for Text Classification.
\newblock In \textit{Advances in Neural Information Processing Systems (NeurIPS)}, pages 649--657, 2015.

\bibitem{greene2006practical}
D.~Greene and P.~Cunningham.
\newblock Practical Solutions to the Problem of Diagonal Dominance in Kernel Document Clustering.
\newblock In \textit{Proceedings of ICML}, pages 377--384, 2006.

\end{thebibliography}

% ============================================
% APPENDIX
% ============================================
\newpage
\appendix
\section{Appendix: Code Listings}

\subsection{Configuration (config.py)}
\begin{lstlisting}[caption={Key configuration parameters}]
MAX_SEQ_LENGTH = 128
TRAIN_BATCH_SIZE = 8
GRADIENT_ACCUMULATION_STEPS = 4
LEARNING_RATE = 2e-5
NUM_EPOCHS = 3
FP16 = True
TFIDF_MAX_FEATURES = 50000
TFIDF_NGRAM_RANGE = (1, 2)
\end{lstlisting}

\subsection{Project Execution Order}
\begin{lstlisting}[caption={Running the complete pipeline}, language=bash]
# Step 1: Prepare datasets
python data_prep.py --dataset all

# Step 2: Exploratory Data Analysis
python eda.py --dataset all

# Step 3: Train baselines
python baselines.py --dataset all

# Step 4: Fine-tune transformers
python train_transformer.py --model all --dataset all

# Step 5: Benchmark comparison
python benchmark.py --dataset all

# Step 6: Launch demo
python app.py
\end{lstlisting}

\end{document}
